\section{Conclusions and Recommendations}
\label{sec:conc}

\subsection{Key Findings}
This research examined carbon markets' potential for supporting rural solar mini-grids in SSA through investigation of commercial viability and integration conditions. Evidence from four developers, five experts and policy officials, and one registry professional across Kenya and Nigeria reveals fundamental structural barriers limiting carbon finance effectiveness.

\textbf{Commercial Viability (RQ1):} Carbon finance functions as supplementary rather than transformative for mini-grid investments. No developer indicated carbon revenue was decisive. For projects under 200~kWp, certification costs represent approximately 10\% of CAPEX (Table \ref{tab:pct-costs}), while fixed auditing costs (100,000--300,000 USD) create minimum scale thresholds individual mini-grids cannot meet (REG-01). Current market pricing of 1--5 USD per \tcooe\ renders transaction costs systematically greater than potential revenues, consistent with transaction cost economics literature.~\cite{michaelowa2005,bellassen}

\textbf{Integration Conditions (RQ2):} Successful integration requires simultaneous convergence across five interdependent domains: technical capacity, financial viability, regulatory frameworks, institutional arrangements, and carbon market design. The atmosfair-MiniGridCo case provides evidence consistent with structural market inadequacy: despite technical success, MiniGridCo has not been replicated across 200+ other Nigerian mini-grids. This reveals that the conditions underlying MiniGridCo, specialized intermediary capacity, absorbed certification costs, patient capital, and registry credibility, are structurally unavailable to typical developers. Post-Paris governance fragmentation~\cite{ahonen-governance} and host government authorization requirements under Article 6 create additional barriers beyond those identified in pre-Paris literature. The February 2026 collapse of KOKO Networks following Kenya's refusal of a Letter of Authorisation~\cite{koko-disrupt} validates the authorization risks this research identified empirically, demonstrating these are not theoretical concerns but active threats to carbon-dependent business models.

\subsection{Implications}
\textbf{For Theory:} This research reveals an additionality paradox: projects with clearest need for additional financing face systematic exclusion due to transaction costs rather than failing additionality criteria.~\cite{michaelowa-evolution} Small mini-grids in remote areas would not proceed without carbon revenue, yet these characteristics create prohibitive certification costs, inverting carbon markets' stated purpose.

The MiniGridCo case illustrates market failure dynamics consistent with institutional economics on coordination problems and public goods provision: aggregation creates collective value but requires institutional capacity that individual actors cannot provide and markets do not produce organically. This confirms carbon market access for small-scale distributed renewable energy requires public sector intervention rather than purely market-based solutions.~\cite{michaelowa2005}

\textbf{For Policy:} Policymakers in Kenya and Nigeria face immediate decisions regarding carbon market integration. Current frameworks operate in silos, with energy regulations~\cite{epra-minigrid-regulations,nerc-minigrid-regulations} and carbon market frameworks~\cite{kenya-climate-change-regulations,ng-carbon-market-policy} lacking coordination. Three critical gaps require attention: (1) clarify carbon revenue treatment in tariff regulations, (2) reform grid arrival compensation to include lost carbon revenues, and (3) distinguish conditional from unconditional NDC targets to avoid disqualifying projects.~\cite{paris}

\textbf{For Practice:} Mini-grid developers should pursue carbon finance only within established aggregation programs, for projects exceeding 500~kWp, or with technical assistance covering certification costs. Independent small-scale developers face structural barriers that individual effort cannot overcome. Alternative mechanisms such as RECs may offer more accessible pathways (REG-01; DEV-KE-02). Carbon registries should develop mini-grid-specific methodologies with proportional MRV requirements and streamline PoAs registration.~\cite{goldStandardFees,verraFees,rrecs}

\subsection{Recommendations}
Successful carbon market integration for mini-grids requires coordinated interventions across multiple domains simultaneously. Priority actions include:

\noindent\textbf{Critical Infrastructure (Public Sector/DFI-Led):}
\begin{enumerate}
    \item Establish publicly-funded aggregation intermediaries in SSA with mandates to absorb upfront certification costs, provide technical support, and coordinate multi-developer PoAs. Market forces have proven insufficient to produce this infrastructure organically.
    \item Subsidize certification costs for aggregated renewable energy projects with registry fee waivers and simplified MRV protocols. DFIs should establish dedicated carbon market access facilities within rural electrification programs.
\end{enumerate}

\noindent\textbf{Regulatory Framework Reforms:}
\begin{enumerate}[resume]
    \item Issue explicit regulatory guidance on carbon revenue treatment within tariff frameworks (e.g., EPRA and the Nigerian Electricity Regulatory Commission, NERC), ensuring environmental benefit streams are not captured through mandatory tariff reductions.
    \item Reform grid arrival compensation frameworks to include projected carbon revenues over remaining crediting periods.~\cite{tenenbaum2014}
    \item Implement tax exemptions for carbon credit revenues to align with renewable energy policy objectives, and establish formal NDC-carbon market coordination distinguishing conditional from unconditional targets.~\cite{glasgow}
\end{enumerate}

\noindent\textbf{Carbon Market Design Adaptations:}
\begin{enumerate}[resume]
    \item Develop mini-grid-specific methodologies acknowledging bundled development co-benefits, proportional MRV requirements for microscale projects (<10,000~\tcooe\ annually), and strong additionality in energy-poor contexts.
    \item Create REC-carbon hybrid instruments providing diversified revenue streams~\cite{rrecs}.
\end{enumerate}

\noindent\textbf{Regional Coordination:}
\begin{enumerate}[resume]
    \item Develop harmonized carbon market frameworks through regional bodies (e.g., the East African Community, EAC, and the Economic Community of West African States, ECOWAS) enabling cross-border aggregation, similar to regional power pool coordination.~\cite{afdb2023}
    \item Integrate carbon finance realistically into national rural electrification master plans as supplementary rather than primary financing mechanism.
\end{enumerate}

\subsection{Future Research}
Three critical research gaps include:
\begin{enumerate}
    \item \textbf{Rigorous impact evaluation:} Quasi-experimental designs with baseline data, matched comparison groups, and longitudinal tracking (3--5 years) would enable robust causal inference on whether carbon finance affects tariffs, service quality, or community outcomes.~\cite{duflo} Key questions: Do carbon-financed projects serve more remote/lower-income communities? Does carbon revenue enable lower tariffs?
    \item \textbf{Financial performance analysis:} Detailed tracking of actual carbon revenue realization rates, timing relative to cash flow needs, audited transaction costs, and revenue volatility would address critical evidence gaps.~\cite{bhattacharyya2016} Understanding which project characteristics enable positive returns would inform targeting of carbon finance mechanisms.
    \item \textbf{Aggregation methods comparison:} Systematic comparison of aggregation models would identify optimal institutional arrangements balancing transaction cost reduction with developer autonomy and local ownership.
\end{enumerate}

\subsection{Conclusion}
This research demonstrates carbon markets face fundamental structural barriers limiting effectiveness for rural mini-grid financing in sub-Saharan Africa. Transaction costs systematically exceed potential revenues for small-scale projects. Methodologies inadequately address mini-grid characteristics. Regulatory frameworks fail to integrate carbon revenue considerations. Most critically, institutional intermediaries necessary for aggregation do not emerge organically from market forces, the atmosfair-MiniGridCo case demonstrates technical feasibility while demonstrating market inadequacy through non-replication.

Carbon market integration requires coordinated public intervention across all five condition domains simultaneously (Table~\ref{tab:rq2}). Addressing single constraints proves insufficient; systemic structural transformation is necessary. This includes publicly-funded aggregation infrastructure, regulatory reforms clarifying carbon revenue treatment, simplified methodologies with proportional MRV requirements, and regional coordination, interventions that markets alone will not produce.~\cite{michaelowa2005}

For the 600 million people in sub-Saharan Africa lacking electricity access,~\cite{iea2023} carbon markets represent one potential supplementary financing mechanism among many necessary tools. Grant finance, concessional loans, innovative business models, and supportive policies all play essential roles.~\cite{esmap2019,tenenbaum2014} Carbon finance can supplement these mechanisms under specific conditions requiring substantial institutional support, but expectations must be calibrated accordingly.

The imperative of universal energy access by 2030~\cite{un2030agenda} demands pragmatic assessment of what works, clear identification of what must change, and sustained commitment to ensuring climate finance mechanisms serve rather than constrain energy access goals. This research indicates that without major structural reforms, particularly publicly-funded aggregation infrastructure and regulatory framework integration, carbon markets will remain largely inaccessible to the mini-grid sector.